\begin{lemma}
\label{lemma-base-change-q-integral-top}
With $f : X \to S$ and $n$ as in Remark \ref{remark-base-change-holds}
assume for some $q \geq 1$ we have $BC(f, n, q - 1)$. Consider
commutative diagrams
$$
\vcenter{
\xymatrix{
X \ar[d]_f &
X' \ar[d]_{f'} \ar[l] &
Y \ar[l]^h \ar[d]^e &
Y' \ar[l]^{\pi'} \ar[d]^{e'} \\
S &
S' \ar[l] &
T \ar[l]_g &
T' \ar[l]_\pi
}
}
\quad\text{and}\quad
\vcenter{
\xymatrix{
X' \ar[d]_{f'} & & Y' \ar[ll]^{h' = h \circ \pi'} \ar[d]^{e'} \\
S' & & T' \ar[ll]_{g' = g \circ \pi}
}
}
$$
where all squares are cartesian, $g$ quasi-compact and quasi-separated, and
$\pi$ is integral surjective. Let $\mathcal{F}$ be an abelian sheaf
on $T_\etale$ annihilated by $n$ and set $\mathcal{F}' = \pi^{-1}\mathcal{F}$.
If the base change map
$$
(f')^{-1}R^qg'_*\mathcal{F}' \longrightarrow R^qh'_*(e')^{-1}\mathcal{F}'
$$
is an isomorphism, then the base change map
$(f')^{-1}R^qg_*\mathcal{F} \to R^qh_*e^{-1}\mathcal{F}$
is an isomorphism.
\end{lemma}

\begin{proof}
Observe that $\mathcal{F} \to \pi_*\pi^{-1}\mathcal{F}'$ is injective
as $\pi$ is surjective (check on stalks). Thus by
Lemma \ref{lemma-base-change-q-injective}
we see that it suffices to show that the base change map
$$
(f')^{-1}R^qg_*\pi_*\mathcal{F}'
\longrightarrow
R^qh_*e^{-1}\pi_*\mathcal{F}'
$$
is an isomorphism. This follows from the assumption because
we have $R^qg_*\pi_*\mathcal{F}' = R^qg'_*\mathcal{F}'$,
we have $e^{-1}\pi_*\mathcal{F}' =\pi'_*(e')^{-1}\mathcal{F}'$, and
we have $R^qh_*\pi'_*(e')^{-1}\mathcal{F}' = R^qh'_*(e')^{-1}\mathcal{F}'$.
This follows from Lemmas
\ref{lemma-integral-pushforward-commutes-with-base-change} and
\ref{lemma-what-integral} and the relative leray spectral sequence
(Cohomology on Sites, Lemma \ref{sites-cohomology-lemma-relative-Leray}).
\end{proof}
